%%%%%%%%%%%%%%%%%%%%%%%%%%%%%%%%%%%%%%%%%%%%%%%%%%%%%%%%%%%%%%%%%%
\section{Conclusions}\label{sec:conclusion}
%%%%%%%%%%%%%%%%%%%%%%%%%%%%%%%%%%%%%%%%%%%%%%%%%%%%%%%%%%%%%%%%%%
% Every \tbd below repeats a \tbd of Sections 4-9 (or the abstract) and is
% filled from the same campaign result; no number here has another source.

\todo{Every qualitative statement below repeats an expected claim of Sections~\ref{sec:verification_exact}--\ref{sec:limitations}; check each against the final results and delete or reverse any that the campaign does not support.}

This paper has presented a quasi-monolithic cell-centred finite volume method for fluid--solid interaction, in which the arbitrary Lagrangian--Eulerian fluid momentum and continuity equations and the total Lagrangian solid momentum equation in velocity form are solved as one nonlinear system by a Jacobian-free Newton--Krylov method, with the fluid mesh motion treated outside the Newton system~\citep{Liu2014, JaimanJoshi2022}.
The method was verified against exact solutions, compared with the partitioned solver of solids4foam on the same discretisation, and assessed for robustness, cost and parallel scaling against Aitken, IQN-ILS and Robin--Neumann partitioned coupling at equal verified accuracy, using the references and the strict coupling convergence test of the companion verification study~\citetodo{companion verification paper}.
The main findings, in the order of the contributions listed in Section~\ref{sec:intro}, are as follows.

\paragraph{Newton operator and preconditioner.}
The compact finite volume linearisations of the fluid, solid and interface-coupling terms are adequate as a preconditioner but not as a Jacobian.
Used as the Newton operator, the assembled approximate Jacobian gives linear convergence with a contraction factor of \tbd{value} per iteration, whereas the matrix-free operator converges in \tbd{number} Newton iterations per step in every case that completes (Section~\ref{sec:tube3d_pc_study}).
With the fluid--solid field-split preconditioner, the Krylov iterations per Newton iteration \tbd{remain below N / grow by factor} over the cases and process counts considered (Sections~\ref{sec:verification_exact} and~\ref{sec:scaling_3d}).
\tbd{One sentence on the preconditioner settings that matter, from the comparison of Section~\ref{sec:tube3d_pc_study}: Schur factorisation, lag, solid block.}

\paragraph{Interface treatment.}
On the ring added-mass problem, the quasi-monolithic coupling reproduces the exact added-mass frequency ratio at order \tbd{order} at every added-mass level, and the amplitude loss of the coupled vibration equals that of the dry ring to within \tbd{difference}, so the coupling adds no measurable damping (Section~\ref{sec:ring_m}).
Replacing the interface pressure gradient taken from the normal momentum equation by a zero normal gradient reduces the order of the frequency ratio to first (observed orders $0.74$--$0.89$) and adds $0.34$ to $2.3$ percentage points of amplitude loss per period (Section~\ref{sec:interface_comparison}); the pressure treatment therefore controls the damping of the coupled vibration as well as its frequency.
Removing the one-step lag of the interface force in the velocity solid, and evaluating the viscous traction from the discrete momentum flux, were each shown during development to remove a loss of accuracy, the latter on a stationary-wall case (Section~\ref{sec:interface_comparison}).
On a deforming wall the traction must also contain the transpose part of the viscous stress, $\mu\bb{\nabla}(\bb{v}\cdot\hat{\bb{n}})$ at the wall: without it the quasi-monolithic wall phase of the pulsatile tube is offset from the partitioned solution by $4.8\times10^{-4}$\,rad on every mesh and time step, and with it, evaluated from the wall velocity, the offset falls tenfold (Sections~\ref{sec:fluid_differences} and~\ref{sec:womersley_consistency}).
Every quantity of the pulsatile tube converges \tbd{at second order} in space and time (observed orders \tbd{range}; Section~\ref{sec:womersley_m}).

\paragraph{Conservative non-conformal coupling.}
Transferring the fluid face forces to the solid with the transpose of the solid-to-fluid velocity transfer balances the total interface force to round-off (relative imbalance below $3\times10^{-15}$) on coarser and finer solid interface meshes, whereas with interpolated traction the imbalance is $3$--$9\times10^{-5}$, the same as on a conformal interface, because the linear-geometry solid integrates the traction over its undeformed face areas (Section~\ref{sec:interface_comparison}).
The choice of transfer changes the frequency ratio of the ring by at most $2.6\times10^{-5}$ relative.
With this transfer, the solid can be refined independently of the fluid, which \tbd{was / was not} used to reduce the solid error on the collapsible channel (Sections~\ref{sec:robustness_collapsible} and~\ref{sec:lim_thin_solid}).

\paragraph{Consistency, robustness and cost.}
On the same meshes and time steps the quasi-monolithic and partitioned solutions differ by \tbd{value} for the flexible-bottom cavity and \tbd{value} for the three-dimensional tube, and the differences decrease under refinement at \tbd{rate} (Sections~\ref{sec:cavity_m} and~\ref{sec:tube3d_consistency}), so the two couplings are consistent to within their discretisation error.
The quasi-monolithic coupling completes \tbd{every cell / which cells} of the robustness maps of the ring and the collapsible channel, including the massless wall, with \tbd{number} Newton iterations per step independently of the density ratio, whereas \tbd{which partitioned couplings} fail above an added-mass ratio of \tbd{value} or below \tbd{steps per period} (Section~\ref{sec:robustness}).
At equal error against the Featflow reference on Hron--Turek FSI3, the quasi-monolithic coupling costs \tbd{factor} times \tbd{more/less} than IQN-ILS, and below an error of \tbd{value} the reference uncertainty, not the coupling, limits the comparison (Section~\ref{sec:fsi3_m}).
In time, the quasi-monolithic coupling converges at order \tbd{order} on the elastic half-cylinder of \citet{Liu2014} (Section~\ref{sec:blob_m}).
On the three-dimensional tube, the method reaches a parallel efficiency of \tbd{value} on \tbd{N} processes (Section~\ref{sec:tube3d_scaling}).
In one case, a thin flag in air, Newton stagnates at the time step at which IQN-ILS converges, and the quasi-monolithic coupling requires \tbd{a quarter} of that step, at \tbd{factor} times the partitioned cost (Section~\ref{sec:lim_foil}): convergence of the method is independent of the added mass but not of the nonlinearity of the coupled problem over a step.

\paragraph{Outlook.}
Several extensions follow from the limitations identified in Section~\ref{sec:limitations}.
First, the solid block of the preconditioner \tbd{is factorised on a single process}, which limits parallel scaling for large solids; an algebraic multigrid solid block with an acceptance test on the converged solution, and a physics-based preconditioner in which a segregated pressure--velocity correction approximates the inverse of the fluid block, are available in the implementation but have not been assessed here.
\todo{Keep the physics-based preconditioner sentence only if Philip wants it named; it was dropped from the paper (Wave 1 decision), but the hook remains in \texttt{newtonIcoFluid} (\texttt{pc\_type physics}).}
Second, a high-order reconstruction for the velocity solids, as already available for the displacement solids, would remove the excess stiffness of thin, bending-dominated structures (Section~\ref{sec:lim_thin_solid}).
Third, the time-step limit of Section~\ref{sec:lim_foil} calls for a globalisation of the Newton iteration beyond a backtracking line search, such as pseudo-transient continuation or adaptive time stepping on Newton failure~\citep{Knoll2004}, or for a convective flux extrapolated from previous time levels (the \texttt{forceImplicitFlux} option of \texttt{newtonIcoFluid}, enabled in every case of this paper), which removes the convective nonlinearity within a step at the price of a time-lagged flux; neither has been assessed.
\todo{Decide (with Section~\ref{sec:lim_foil}) whether pseudo-transient continuation is named as future work.}
Finally, larger three-dimensional cases than the tube of Section~\ref{sec:scaling_3d} are needed to establish the cost of the method at the scale at which partitioned finite volume FSI solvers are applied, and, since the fluid mesh motion is already outside the Newton system, the coupled residual could be combined with overset meshes~\citep{miller2014overset} for motions that a body-fitted ALE mesh cannot follow.
